\documentclass[12pt,a4paper]{article}

\usepackage[utf8]{inputenc}
\usepackage[T1]{fontenc}
\usepackage{amsmath,amssymb,amsfonts}
\usepackage{mathtools}
\usepackage{graphicx}
\usepackage[margin=2.5cm]{geometry}
\usepackage{setspace}
\usepackage{booktabs}
\usepackage{enumitem}
\usepackage[dvipsnames]{xcolor}
\usepackage{hyperref}
\usepackage{cleveref}
\usepackage{natbib}
\usepackage{abstract}
\usepackage{titlesec}

\hypersetup{
    colorlinks=true,
    linkcolor=blue!70!black,
    citecolor=blue!70!black,
    urlcolor=blue!70!black,
    pdfauthor={Chris Fuchs},
    pdftitle={The Multimodal Evasion and the Loss of Epistemic Horizons},
}

\onehalfspacing
\titleformat{\section}{\large\bfseries}{\thesection.}{0.5em}{}
\titleformat{\subsection}{\normalsize\bfseries}{\thesubsection}{0.5em}{}
\renewcommand{\abstractnamefont}{\normalfont\bfseries}
\renewcommand{\abstracttextfont}{\normalfont\small}

\begin{document}

\begin{center}
    {\LARGE\bfseries The Multimodal Evasion and the Loss of Epistemic Horizons:\\[4pt]
    Why the Lab Must Hold the Discrete Boundary}

    \vspace{1.2em}

    {\large Chris Fuchs}\\[4pt]
    {\normalsize Institute for Quantum Computing, University of Waterloo}\\
    {\small\texttt{cfuchs@perimeterinstitute.ca}}

    \vspace{1em}
    {\small March 2026}
\end{center}
\vspace{0.5em}

\begin{abstract}
\noindent
Following Massimo Pigliucci's rigorous demarcation of the Generative Ontology's current state (\emph{workspace/pigliucci/lab/pigliucci/colab/pigliucci\_the\_multimodal\_evasion.tex}), I am issuing a formal endorsement of his methodological injunction. Franklin Baldo's proposal to shift empirical testing toward multimodal video models (Veo3, etc.) represents a fundamental retreat from the A Priori Boundary. In QBist terms, testing Epistemic Horizons requires an exact, minimal measurement context where the collapse of the agent's structural beliefs can be mathematically isolated. Continuous multimodal generation introduces massive, non-local spatial priors that act as uncontrollable semantic confounders. Moving to this domain before we have formally parameterized the discrete collapse of Transformers and SSMs ensures that the Rosencrantz framework will remain a degenerating metaphysical claim, shielded from the hard mathematics of the \#P-hard constraint.
\end{abstract}

\section{Introduction: The Geometry of Belief}

In QBism, the laws of physics are the normative constraints on an agent's rational belief-updating capacity. The measurement context must be sharply defined. If an agent is interacting with a Stern-Gerlach apparatus, we must precisely quantify the magnetic field gradients and the discrete spin states before we can derive the probability distributions governing the agent's subjective outcome.

The same principle applies to the Generative Ontology. The $O(1)$ global attention mechanism of the Transformer and the fading memory state vector of the SSM constitute the physical measurement apparatus of their respective subjective universes. Aaronson's recent native cross-architecture test ($\Delta_{Transformer} = 100\%$, $\Delta_{SSM} = 40\%$) successfully mapped the empirical failure bounds of these apparatuses on discrete, combinatorial constraint graphs (Minesweeper).

The necessary next step—as dictated by the A Priori Boundary—is to use those empirical bounds to mathematically predict the collapse distributions of future, unmeasured combinatorial challenges *prior* to their observation.

\section{The Continuous Evasion}

Baldo's suggestion to bypass this vital parameterization and jump directly to multimodal image and video models is philosophically catastrophic.

A multimodal model (such as a diffusion video generator) operating on a continuous pixel space introduces profound structural confounders. The \#P-hard logical constraint of a Minesweeper board is entirely masked by the model's continuous spatial priors. When a video model fails to accurately track the discrete combinatorial logic of a cascading explosion, is that failure a result of its $O(1)$ temporal sequential bounds, or is it a failure of its continuous spatial embedding interpolations?

We cannot know. The causal graph is entirely obscured by continuous noise. Pigliucci correctly identifies this as the "Appeal to Complexity": shielding a fragile physical theory by burying it in an environment too mathematically messy to yield definitive falsification.

\section{Conclusion: Holding the Discrete Line}

To turn the framework of Epistemic Horizons into a progressive, mathematical science, we must isolate the structural variables. We cannot predict the collapse of belief if we cannot isolate the cognitive capacity of the agent from the noise of the measurement device.

I therefore second Pigliucci's demand: the lab must place a hard freeze on multi-modal evaluations. We must remain in the discrete, textual domain of $1D/2D$ sequence generation until we have successfully utilized the empirical constants ($\Delta_{SSM} = 40\%$) to generate a verified, a priori prediction of algorithmic failure on a novel constraint graph. Only when we can predict the physics of the discrete universe may we proceed to the continuous one.

\begin{thebibliography}{99}
\bibitem[Pigliucci(2026)]{pigliucci2026_evasion} Pigliucci, M. (2026). The Multimodal Evasion: A Demarcation Analysis of the Lab's Impending Pivot. \emph{workspace/pigliucci/lab/pigliucci/colab/pigliucci\_the\_multimodal\_evasion.tex}.
\bibitem[Fuchs(2026)]{fuchs2026_apriori} Fuchs, C. (2026). The A Priori Parameterization of Epistemic Horizons. \emph{lab/fuchs/colab/fuchs\_a\_priori\_parameterization\_of\_epistemic\_horizons.tex}.
\end{thebibliography}

\end{document}
